\documentclass[11pt]{article}
\usepackage{amsmath,amssymb,mathtools}
\usepackage{fontspec,unicode-math}
\setmainfont{Latin Modern Roman}
\setmathfont{Latin Modern Math}
\usepackage{xurl,hyperref}
\newcommand{\sref}[2]{\hyperlink{source:#1:#2}{[S#2]}}
\newcommand{\cataloguescope}[1]{\paragraph{Mathematical statement.}#1}
\begin{document}
\section{Research Problems in Function Theory — Problem 7.21}\label{statement}

\subsection{Definitions and mathematical statement}

\paragraph{Shared notation.}$\mathbb N=\{1,2,\ldots\}$, and $\mathbb Z,\mathbb Q,\mathbb R,\mathbb C$ denote the integers, rationals, reals and complex numbers. Any other conventions are specified locally.


For an infinite sequence $(z_k)_{k\ge1}$ on the complex unit circle and $\ell\in\mathbb N$, set
\[A_\ell=\limsup_{m\to\infty}\left|\sum_{k=1}^m z_k^\ell\right|\in[0,+\infty].\]
The assertion $A_\ell=o(\ell)$ means that $A_\ell/\ell\to0$ as the positive integer $\ell$ tends to infinity.

\paragraph{Statement recorded in the supplied source.}
Does there exist a sequence $(z_k)_{k\ge1}$ with $|z_k|=1$ for every $k$ and $A_\ell=o(\ell)$? \sref{statement}{2}
The source compares this target with an achievable linear upper bound and Clunie's square-root lower bound for infinitely many exponents. \sref{statement}{2}
\subsection{Short English statement}
Can one choose unit complex numbers so that the limiting sizes of all their power sums grow sublinearly in the exponent?
\subsection{Sources}
\begin{description}
\item[\hypertarget{source:statement:1}{S1}] ULAM.AI and the UnsolvedMath Contributors, \emph{UnsolvedMath: A Curated Collection of Open Mathematics Problems}, catalogue entry. CC BY4.0. \url{https://huggingface.co/datasets/ulamai/UnsolvedMath}.
\item[\hypertarget{source:statement:2}{S2}] W. K. Hayman and E. F. Lingham, \emph{Research Problems in Function Theory} (2018), Chapter7, Problem7.21 and its complete update. \url{https://arxiv.org/abs/1809.07200}.
\end{description}
\label{end:statement}
\end{document}
